@@M0@@
@@M1@@
And to establish II, it suffices to see that, after shrinking @@M2@@, @@M3@@, @@M4@@ and @@M5@@ are quasi-affine.

Consider @@M6@@ first. Replacing @@M7@@ by @@M8@@ and @@M9@@ by @@M10@@ reduces us to the case where @@M11@@ is integral over @@M12@@. Since @@M13@@ is surjective, one of the @@M14@@, say @@M15@@, dominates @@M16@@, and we have the two sequences
@@M17@@
@@M18@@
Thus, knowing that the @@M19@@, and therefore also @@M20@@, are representable and that the composite morphism is quasi-affine, we conclude that @@M21@@ is quasi-affine. In view of the proof of 1.2, for @@M22@@ we are therefore reduced to proving II'. \uline{Theorem} 1.1 \uline{when} @@M23@@ \uline{and} @@M24@@ \uline{are integral over} @@M25@@.

For @@M26@@ and @@M27@@, we must treat the case of a morphism @@M28@@ where @@M29@@ is reduced, the @@M30@@ are the irreducible components of @@M31@@ and are integral over @@M32@@. Now we may suppose that @@M33@@ has connected fibers. Indeed, let @@M34@@ be the geometric point of @@M35@@. Then the irreducible components of @@M36@@ are the @@M37@@ and, after shrinking @@M38@@, we may suppose that @@M39@@ if and only if

